% エントリ 22(2026-11-14): BH の年齢の分布の模型と、年齢について周辺化した窓
\section{[E22, 2026-11-14] BH の年齢の分布の模型と、年齢について周辺化した窓}\label{sec:e22}

固定点の領域では、スピンは $\tau=\mu(f/\Mpl)^2T$ だけで決まる(E12、DESIGN.md D9)。$f$ の制限は年齢 $T$ の仮定にそのまま比例して動くのに、これまで年齢は一つの値に固定していた(2 世代目は $3\times10^7$--$3\times10^8\yr$、1 世代目は $6\times10^9\yr$)。集団解析では年齢を分布として周辺化する。そのための模型を、解析が呼べる形で書いた(ユーザー指示「年齢の模型を書く」、2026-11-14)。

結論を先に書く。
\begin{enumerate}
\item \textbf{模型}: 1 世代目の年齢は「遅れの分布 $\times$ 宇宙の星形成率」、2 世代目の年齢は「時間の分布の形 $\times$ 親の合体率」で書く(式~\eqref{eq:e22p1}、\eqref{eq:e22p2})。合体の赤方偏移ごとに、宇宙の年齢で上限が切れる。部品は \texttt{share/ages/ages.py}(numpy だけ)で、年齢の分位点を求積の節として返す。
\item \textbf{1 世代目}: 孤立連星の標準の形(遅れ $\propto1/T$、10 Myr 以上)に星形成率の重みをつけると、$z=0.2$ で合体する BH の年齢の中央値は 1.5 Gyr で、16--84\% は $5.5\times10^7$--$6.9\times10^9\yr$ に広がる。22\% は $10^8\yr$ より若い。遅れを $T$ について一様にすると中央値 6.6 Gyr になり、これまで使ってきた 6 Gyr(古い球状星団)にあたる。
\item \textbf{2 世代目}: 年齢の分布を示した文献はない(E9)。形をパラメータで書き、目安として四つの組(広い、星団の目安、遅い、AGN 円盤)を置いた。中央値は $3\times10^5\yr$ から $3\times10^9\yr$ まで 4 桁にわたる。
\item \textbf{窓の条件は年齢の比で書ける}: 質量 $M$ の 1 世代目が回転を失い、GW241011 の主星が回転を保つ $f$ があるのは、$T_1/T_2\ge R(M)$ のときである。$R$ は $f$ によらず、$10$、15、20、$25\Msun$ で 299、36、7.4、1.9(生まれたスピン 0.78、$\delta=0.05$)。
\item \textbf{年齢について周辺化すると、窓は「ある・ない」ではなく程度になる}。同じ $f$ で「2 世代目が保つ確率」と「1 世代目が失う確率」の小さいほうを最大にした値 $W$ は、星団の目安の 2 世代目と古い 1 世代目で、$20\Msun$ が 0.88、$15\Msun$ が 0.71、$10\Msun$ が 0.33。1 世代目を既定の広い年齢にすると 0.60、0.47、0.26 に下がる。2 世代目が遅い($3\times10^8\yr$ 以上)なら、$20\Msun$ でも 0.24 以下である。
\item \textbf{$f$ を測る物差しは 1 世代目の年齢である}。年齢をすべて $k$ 倍にして $f$ を $k^{-1/2}$ 倍にしても何も変わらないので、2 世代目の年齢の尺度は、$f$ と一緒にデータから制限される量になる。
\end{enumerate}

計算は \texttt{calc/e22\_ages.py}(出力 \texttt{calc/out/e22.json})。$\chi_\mathrm{SR}$ は単体版(E21 の版)、$\mu=0.71\times10^{-12}\eV$。

\subsection{定義と模型}

年齢 $T$ は、BH が生まれてから、その BH を含む連星が合体するまでの時間である。超放射はこのあいだ働く。合体の赤方偏移を $z_m$、そのときの宇宙の年齢を $t_m$ と書く。宇宙論は平坦な $\Lambda$CDM で、$H_0=67.4\,\mathrm{km\,s^{-1}Mpc^{-1}}$、$\Omega_m=0.315$(\arxiv{1807.06209} の要旨の値)、$t(z)=\frac{2}{3H_0\sqrt{\Omega_\Lambda}}\,\mathrm{arcsinh}\bigl[\sqrt{\Omega_\Lambda/\Omega_m}\,(1+z)^{-3/2}\bigr]$ である(いまの宇宙の年齢は 13.80 Gyr、$z=0.2$ で 11.28 Gyr)。

\paragraph{1 世代目}
$T$ は星形成から合体までの遅れとする(BH は星形成の 3--20 Myr あとに生まれるが、その差は無視する)。時刻 $t_m$ に合体する BH のうち、年齢が $T$ のものは、時刻 $t_m-T$ に生まれて、遅れ $T$ で合体したものである。したがって
\begin{equation}
p_1(T\,|\,z_m)\propto T^{-\kappa_1}\,\psi(t_m-T),\qquad T_\mathrm{min,1}\le T\le t_m-t_\mathrm{start}.
\label{eq:e22p1}
\end{equation}
$T^{-\kappa_1}$ は遅れの分布、$\psi$ は宇宙の星形成率で、Madau \& Dickinson(\arxiv{1403.0007})の式 15
\begin{equation}
\psi(z)\propto\frac{(1+z)^{2.7}}{1+[(1+z)/2.9]^{5.6}}
\end{equation}
を使う($z=1.9$ に山がある)。$t_\mathrm{start}$ は星形成の始まりで、$z=20$(0.18 Gyr)にとる。

\paragraph{2 世代目}
$T$ は、親の合体(この BH ができた時刻)から次の合体までの時間とする。同じ考え方で
\begin{equation}
p_2(T\,|\,z_m)\propto p_D(T)\,R_1(t_m-T),\qquad T_\mathrm{lo}\le T\le\min(T_\mathrm{hi},\,t_m-t_\mathrm{start}-T_\mathrm{min,1}),
\label{eq:e22p2}
\end{equation}
$p_D$ は時間の分布の形、$R_1(t)=\int T'^{-\kappa_1}\psi(t-T')\,dT'$ は 1 世代目どうしの合体率(式~\eqref{eq:e22p1} と同じ模型から出す)である。$p_D$ は対数一様($\propto1/T$)か、対数正規($\ln T$ が中央値 $T_\mathrm{med}$、幅 $s$ dex の正規分布)を選ぶ。

\begin{table}[H]
\centering\footnotesize
\begin{tabular}{llll}
\toprule
パラメータ & 既定 & 意味 & 根拠\\
\midrule
$\kappa_1$ & 1 & 1 世代目の遅れのべき & 孤立連星の標準の形 $p(\tau)\propto\tau^{-1}$(\arxiv{2105.06491} Sec.~1)\\
$T_\mathrm{min,1}$ & $10^7\yr$ & 1 世代目の遅れの下限 & 孤立連星の予言(同 Sec.~4)\\
$p_D$ & 対数一様 & 2 世代目の時間の分布の形 & 文献に分布がないので、いちばん仮定の少ない形\\
$T_\mathrm{lo}$、$T_\mathrm{hi}$ & $10^6$、$10^{10}\yr$ & 2 世代目の年齢の範囲 & 星団の模型の目安の下端から、宇宙の年齢まで\\
$T_\mathrm{med}$、$s$ & $10^8\yr$、1 dex & 対数正規のときの中央値と幅 & (形を変えるときに使う)\\
\bottomrule
\end{tabular}
\caption{年齢の模型のパラメータ(\texttt{ages.AgeModel} の引数)。星形成率の重み $\psi$ と親の合体率の重み $R_1$ は、それぞれ外すこともできる。}
\label{tab:e22par}
\end{table}

\paragraph{文献にあること}
抜き書きは \texttt{Refs/notes/astro\_inputs.md} の 3 節(2026-11-03 に本文で照合)。
\begin{itemize}
\item \textbf{1 世代目}: 孤立連星は $p(\tau)\propto\tau^{-1}$、$\tau_\mathrm{min}\sim10$ Myr が標準の予言で、GWTC-2 からの推定は $\tau_{50\%}=2.8^{+3.3}_{-2.6}$ Gyr(\arxiv{2105.06491})。古い球状星団から低赤方偏移で合体するものは $\gtrsim6$ Gyr(\arxiv{1911.00018} Fig.~15)。
\item \textbf{2 世代目}: 年齢の分布を直接示した論文は見つからなかった。近いものは、(i) 遅れの合計の最小値の目安が $(50+50)\Msun$ で $5\times10^7\yr$、$(10+10)\Msun$ で $3\times10^8\yr$(\arxiv{2004.00650} Fig.~2 から読んだ値)、(ii) 星団の模型で次の世代ができるまでの時間が 0.3 Myr から 7 Gyr、山が 10--100 Myr(\arxiv{2103.05016} Figs.~1--2 から読んだ値。増分の時間と解釈した)、(iii) AGN 円盤では円盤の寿命 0.1--10 Myr のうちに全世代が合体する(\arxiv{2311.18548}、\arxiv{2410.16515})。
\item \textbf{先行研究が置いた仮定}: \arxiv{2609.02678} はどの BH も 1 Myr--2.2 Gyr の対数一様、\arxiv{2510.26931} は GW241011 の主星を控えめに $10^5\yr$、\arxiv{1703.06223} は $10^{-4}$ Gyr から宇宙の年齢までの対数一様。
\end{itemize}

\paragraph{比べるための組}
パラメータの値を変えた組を、1 世代目に三つ、2 世代目に四つ置いた(表~\ref{tab:e22q})。2 世代目の組は上の目安に合わせたもので、文献にある分布そのものではない。
\begin{itemize}
\item 1 世代目: 「既定」、「古い」($\kappa_1=0$。遅れが $T$ について一様)、「重みなし」(星形成率の重みを外す。対数一様)。
\item 2 世代目: 「広い」(既定)、「星団の目安」(対数正規、中央値 $5\times10^7\yr$、幅 0.7 dex)、「遅い」($T_\mathrm{lo}=3\times10^8\yr$)、「AGN 円盤」($10^4$--$10^7\yr$ の対数一様)。
\end{itemize}

\subsection{部品と確かめ}

\texttt{ages.py} は、合体の時刻 160 点 $\times$ $\ln T$ 800 点の表で累積分布を作り、分位点を補間で返す。作るのに 0.04 秒、$2\times10^5$ 個の赤方偏移で 32 個の節を出すのに 0.4 秒である。確かめたことは次のとおり(\texttt{e22\_ages.py check})。
\begin{itemize}
\item 宇宙の年齢は、$\int dz/[(1+z)H(z)]$ の数値積分と $3\times10^{-10}$ で一致する。
\item 重みを外したときの分位点は、解析的な値(対数一様、一様、対数正規)と 0.5\% 以内で一致する。
\item 重みつきの累積分布は、表を使わずに別に書いた数値積分と $3\times10^{-4}$ 以内で一致する(四つの組 $\times$ 赤方偏移 5 点 $\times$ 確率 5 点)。分位点と累積分布の往復の差は $4\times10^{-4}$ 以下。
\item 40 万個の標本の経験分布と累積分布の差は 0.002 以下で、標本の数から期待される大きさである。
\item 星形成の始まりを $z=10$ にしても、年齢の中央値は 0.1\% しか変わらない。
\end{itemize}

\subsection{年齢の分布}

\begin{table}[H]
\centering\small
\begin{tabular}{llccccc}
\toprule
世代 & 組 & 5\% & 16\% & 50\% & 84\% & 95\%\\
\midrule
1 & 既定 & $1.7\times10^7$ & $5.5\times10^7$ & $1.5\times10^9$ & $6.9\times10^9$ & $8.6\times10^9$\\
1 & 古い & $1.4\times10^9$ & $3.4\times10^9$ & $6.6\times10^9$ & $8.6\times10^9$ & $9.6\times10^9$\\
1 & 重みなし & $1.4\times10^7$ & $3.1\times10^7$ & $3.3\times10^8$ & $3.6\times10^9$ & $7.8\times10^9$\\
\midrule
2 & 広い & $1.8\times10^6$ & $6.0\times10^6$ & $2.7\times10^8$ & $5.3\times10^9$ & $8.4\times10^9$\\
2 & 星団の目安 & $3.6\times10^6$ & $1.0\times10^7$ & $5.2\times10^7$ & $2.7\times10^8$ & $8.5\times10^8$\\
2 & 遅い & $3.9\times10^8$ & $6.9\times10^8$ & $3.1\times10^9$ & $7.6\times10^9$ & $9.2\times10^9$\\
2 & AGN 円盤 & $1.4\times10^4$ & $3.0\times10^4$ & $3.2\times10^5$ & $3.3\times10^6$ & $7.1\times10^6$\\
\bottomrule
\end{tabular}
\caption{年齢の分位点 [yr]。1 世代目は $z_m=0.2$、2 世代目は $z_m=0.05$(GW241011 の赤方偏移)。}
\label{tab:e22q}
\end{table}

1 世代目の既定の組は広い。$z_m=0.2$ で、$3\times10^7\yr$ より若いものが 10\%、$10^8\yr$ より若いものが 22\%、$10^9\yr$ より若いものが 45\% ある。星形成率の重みは年齢を長い側に寄せる(中央値は、重みなしの $3.3\times10^8\yr$ から $1.5\times10^9\yr$ になる)。昔のほうが星形成率が高く、長い遅れのものが多く生まれているためである。中央値 1.5 Gyr は、GWTC-2 からの推定 $2.8^{+3.3}_{-2.6}$ Gyr の幅の中にある。$\kappa_1=0$ の「古い」組は中央値 6.6 Gyr で、E10--E13 が 1 世代目の年齢に使った 6 Gyr にあたる。合体が遠いほど年齢は短くなる(既定の中央値は、$z_m=0.05$、0.5、1 で 2.5、0.70、0.29 Gyr)。

2 世代目の「広い」組では、親の合体率の重みで中央値が $1.0\times10^8\yr$ から $2.7\times10^8\yr$ に延びる。$5\times10^7\yr$ より若いものは 35\%、$3\times10^8\yr$ より若いものは 51\% である。「星団の目安」ではそれぞれ 49\%、85\%。

\begin{figure}[H]
\centering
% E22: (a) 年齢の分布、(b) 年齢について周辺化した、回転を保つ確率と失う確率。データは calc/e22_ages.py
\begin{tikzpicture}
\begin{groupplot}[group style={group size=2 by 1, horizontal sep=17mm}, width=0.48\linewidth, height=0.44\linewidth, note, xmode=log]
\nextgroupplot[xmin=1e4, xmax=2e10, ymin=0, ymax=3.4, xlabel={年齢 $T$ [年]}, ylabel={$dP/d\log_{10}T$},
  title={\small (a) 年齢の分布},
  legend style={font=\scriptsize, at={(0.03,0.97)}, anchor=north west, fill=white, fill opacity=0.85, text opacity=1, draw=none},
  legend cell align=left]
  \addplot[cA, line width=1.0pt] table[x=T, y=g1_field] {figs/data/e22_pdf.dat}; \addlegendentry{1 世代目、既定}
  \addplot[cA, line width=1.0pt, dashed] table[x=T, y=g1_old] {figs/data/e22_pdf.dat}; \addlegendentry{1 世代目、古い}
  \addplot[cB, line width=1.0pt] table[x=T, y=g2_wide] {figs/data/e22_pdf.dat}; \addlegendentry{2 世代目、広い}
  \addplot[cB, line width=1.0pt, dashed] table[x=T, y=g2_peak] {figs/data/e22_pdf.dat}; \addlegendentry{2 世代目、星団の目安}
  \addplot[cB, line width=1.0pt, densely dotted] table[x=T, y=g2_slow] {figs/data/e22_pdf.dat}; \addlegendentry{2 世代目、遅い}
  \addplot[cM, line width=1.0pt] table[x=T, y=g2_agn] {figs/data/e22_pdf.dat}; \addlegendentry{2 世代目、AGN 円盤}
\nextgroupplot[xmin=3e10, xmax=1e14, ymin=0, ymax=1.05, xlabel={$f$ [GeV]}, ylabel={確率},
  title={\small (b) 年齢について周辺化した確率},
  legend style={font=\scriptsize, at={(0.5,-0.30)}, anchor=north, fill=white, draw=none}, legend columns=2,
  legend cell align=left]
  \addplot[cB, line width=1.0pt] table[x=f, y=wide_078] {figs/data/e22_prob_keep.dat}; \addlegendentry{保つ: 2 世代目、広い}
  \addplot[cA, line width=1.0pt] table[x=f, y=field_20] {figs/data/e22_prob_drop.dat}; \addlegendentry{失う: 1 世代目 $20\,M_\odot$、既定}
  \addplot[cB, line width=1.0pt, dashed] table[x=f, y=peak_078] {figs/data/e22_prob_keep.dat}; \addlegendentry{保つ: 2 世代目、星団の目安}
  \addplot[cA, line width=1.0pt, dashed] table[x=f, y=old_20] {figs/data/e22_prob_drop.dat}; \addlegendentry{失う: 1 世代目 $20\,M_\odot$、古い}
  \addplot[cB, line width=1.0pt, densely dotted] table[x=f, y=slow_078] {figs/data/e22_prob_keep.dat}; \addlegendentry{保つ: 2 世代目、遅い}
  \addplot[cC, line width=1.0pt, dashed] table[x=f, y=old_10] {figs/data/e22_prob_drop.dat}; \addlegendentry{失う: 1 世代目 $10\,M_\odot$、古い}
\end{groupplot}
\end{tikzpicture}

\caption{(a) 年齢の分布(1 世代目は $z_m=0.2$、2 世代目は $z_m=0.05$)。(b) 年齢について周辺化した確率($\mu=0.71\times10^{-12}\eV$)。「保つ」は、$19.6\Msun$ で 0.78 のスピンで生まれた 2 世代目(GW241011 の主星)が 0.69 以上にとどまる確率、「失う」は、0.7 で生まれた 1 世代目が 211 の臨界 $+0.05$ より下にいる確率。}
\label{fig:e22}
\end{figure}

\subsection{窓の条件を年齢の比で書く}

E10 からの説明用の判定をそのまま使う。GW241011 の主星($19.6\Msun$、生まれたスピン $\chi_b$)が年齢 $T_2$ のあいだ 0.69 を下回らない $f$ の上限を $f_\mathrm{max}$、質量 $M$ の 1 世代目(生まれたスピン 0.7)が年齢 $T_1$ のうちに「211 の臨界 $+\delta$」まで落ちる $f$ の下限を $f_\mathrm{low}(M)$ とする。固定点の領域では $f_\mathrm{max}\propto T_2^{-1/2}$、$f_\mathrm{low}\propto T_1^{-1/2}$ である(数値でも $f\sqrt T$ は $10^{-4}$ まで一定)。窓 $f_\mathrm{low}(M)\le f_\mathrm{max}$ は
\begin{equation}
\frac{T_1}{T_2}\ \ge\ R(M)\equiv\left(\frac{f_\mathrm{low}(M)\sqrt{T_1}}{f_\mathrm{max}\sqrt{T_2}}\right)^2
\label{eq:e22R}
\end{equation}
と書ける。$R(M)$ は $f$ にも年齢にもよらず、質量と、生まれたスピン、判定の幅 $\delta$ だけで決まる(表~\ref{tab:e22R})。

\begin{table}[H]
\centering\small
\begin{tabular}{llcccccc}
\toprule
$\chi_b$ & $\delta$ & $8\Msun$ & 10 & 12 & 15 & 20 & 25\\
\midrule
0.78 & 0.05 & 934 & 299 & 116 & 36 & 7.4 & 1.9\\
0.78 & 0.02 & 1650 & 522 & 203 & 62 & 13 & 3.6\\
0.72 & 0.05 & 2560 & 820 & 320 & 99 & 20 & 5.3\\
0.72 & 0.02 & 4520 & 1430 & 557 & 171 & 36 & 10\\
\bottomrule
\end{tabular}
\caption{窓が開くのに要る年齢の比 $R(M)=T_1/T_2$ の下限($\mu=0.71\times10^{-12}\eV$)。}
\label{tab:e22R}
\end{table}

これは E19 の窓の表と同じ内容である。$T_1=6\times10^9\yr$、$T_2=5\times10^7\yr$(比 120)なら $11.9\Msun$ 以上($\chi_b=0.78$、$\delta=0.05$)で、E19 の値と一致する。年齢を、組の中央値に固定して読むと次のようになる($\delta=0.05$、$\chi_b=0.78$--0.72)。
\begin{center}\small
\begin{tabular}{llccc}
\toprule
1 世代目 & 2 世代目 & 中央値の比 & 回転を失える 1 世代目の最小質量\\
\midrule
古い(6.6 Gyr) & 星団の目安($5.2\times10^7\yr$) & 127 & 11.8--14.3$\Msun$\\
古い & 広い($2.7\times10^8\yr$) & 25 & 16.1--19.3$\Msun$\\
既定(1.5 Gyr) & 星団の目安 & 29 & 15.6--18.8$\Msun$\\
既定 & 広い & 5.7 & 21.0--24.8$\Msun$\\
古い & 遅い($3.1\times10^9\yr$) & 2.1 & 24.6--28.5$\Msun$\\
\bottomrule
\end{tabular}
\end{center}
窓が開くかどうかは、年齢の比で 2 桁動く(最後の行の上端は、$(211,322)$ の固定点で印 4 になる質量($\alpha\gtrsim0.15$、$28\Msun$ 以上)にかかるので粗い)。

\subsection{年齢について周辺化した窓}

年齢を分布にすると、同じ $f$ でも、若い BH は回転を保ち、古い BH は失う。年齢の節 400 個で $\chi_\mathrm{SR}$ を引き、次の二つの確率を $f$ の関数として出した(図~\ref{fig:e22}(b))。
\begin{itemize}
\item $P_\mathrm{keep}(f)$: GW241011 の主星($19.6\Msun$、$z_m=0.05$、生まれたスピン 0.72 か 0.78)が 0.69 以上にとどまる確率。
\item $P_\mathrm{drop}(f;M)$: 質量 $M$ の 1 世代目($z_m=0.2$、生まれたスピン 0.7)が「211 の臨界 $+0.05$」より下にいる確率。
\end{itemize}
0.05 は説明用の幅で、集団解析では使わない(D9)。印 3・4・5 の点は 3.5\% 以下である($30\Msun$ は、$(211,322)$ の固定点で 411・433・544 が育ちうる印 4 が 9 割を越えるので、使わなかった)。

\begin{table}[H]
\centering\footnotesize
\begin{tabular}{lccc@{\qquad}lccc}
\toprule
\multicolumn{4}{c}{$P_\mathrm{keep}$ が 0.9、0.5、0.1 になる $f$ [GeV]} & \multicolumn{4}{c}{$P_\mathrm{drop}$ が 0.1、0.5、0.9 になる $f$ [GeV]}\\
2 世代目($\chi_b=0.78$) & 0.9 & 0.5 & 0.1 & 1 世代目 & 0.1 & 0.5 & 0.9\\
\midrule
広い & $1.6\times10^{11}$ & $8.2\times10^{11}$ & $7.7\times10^{12}$ & 既定、$10\Msun$ & $2.6\times10^{12}$ & $6.0\times10^{12}$ & $4.3\times10^{13}$\\
星団の目安 & $6.4\times10^{11}$ & $1.9\times10^{12}$ & $5.3\times10^{12}$ & 既定、$15\Msun$ & $9.1\times10^{11}$ & $2.1\times10^{12}$ & $1.5\times10^{13}$\\
遅い & $1.5\times10^{11}$ & $2.4\times10^{11}$ & $6.0\times10^{11}$ & 既定、$20\Msun$ & $4.1\times10^{11}$ & $9.4\times10^{11}$ & $6.8\times10^{12}$\\
AGN 円盤 & $6.0\times10^{12}$ & $2.4\times10^{13}$ & $9.5\times10^{13}$ & 古い、$10\Msun$ & $2.4\times10^{12}$ & $2.9\times10^{12}$ & $4.8\times10^{12}$\\
 & & & & 古い、$15\Msun$ & $8.5\times10^{11}$ & $9.9\times10^{11}$ & $1.6\times10^{12}$\\
 & & & & 古い、$20\Msun$ & $3.8\times10^{11}$ & $4.5\times10^{11}$ & $7.5\times10^{11}$\\
\bottomrule
\end{tabular}
\caption{確率が決まった値を横切る $f$。$\chi_b=0.72$ では、$P_\mathrm{keep}$ の $f$ は 0.60 倍になる。}
\label{tab:e22f}
\end{table}

年齢が 4 桁に広がる組(2 世代目の「広い」、1 世代目の「既定」)では、確率が 0.9 から 0.1 に変わるのに $f$ で 1.2--1.7 桁かかる。年齢の幅が狭い「古い」1 世代目では 0.3 桁である。$f=10^{12}\GeV$ での $20\Msun$ の 1 世代目のスピンは、既定の組で 16--84\% が 0.41--0.67(中央値 0.45)に散らばり、古い組では 0.40--0.42 に集まる。積み上がりの幅は、年齢の分布で決まる。

\paragraph{同じ $f$ で両立する度合い}
窓を一つの数にまとめる。
\begin{equation}
W(M)\equiv\max_f\,\min\bigl[P_\mathrm{keep}(f),\,P_\mathrm{drop}(f;M)\bigr].
\end{equation}
年齢を一つの値に固定すれば、窓があるとき 1、ないとき 0 になる量である。

\begin{table}[H]
\centering\small
\begin{tabular}{lcccc@{\qquad}cccc}
\toprule
 & \multicolumn{4}{c}{1 世代目: 既定} & \multicolumn{4}{c}{1 世代目: 古い}\\
2 世代目 & $10\Msun$ & 15 & 20 & 25 & $10\Msun$ & 15 & 20 & 25\\
\midrule
広い & 0.26 & 0.39 & 0.48 & 0.55 & 0.29 & 0.47 & 0.59 & 0.71\\
星団の目安 & 0.26 & 0.47 & 0.60 & 0.69 & 0.33 & 0.71 & 0.88 & 0.94\\
遅い & 0.00 & 0.00 & 0.19 & 0.36 & 0.00 & 0.00 & 0.24 & 0.51\\
AGN 円盤 & 0.67 & 0.79 & 0.89 & 0.96 & 0.93 & 0.99 & 1.00 & 1.00\\
\bottomrule
\end{tabular}
\caption{同じ $f$ で両立する度合い $W(M)$。GW241011 の主星の生まれたスピンは 0.78(0.72 では最大 0.22 小さい。\texttt{e22.json})。}
\label{tab:e22W}
\end{table}

読み方は次のとおりである。
\begin{itemize}
\item \textbf{いちばん有利な組でも、軽い 1 世代目は足りない}。星団の目安の 2 世代目と古い 1 世代目で、$20\Msun$ は 0.88、$15\Msun$ は 0.71 だが、$10\Msun$ は 0.33 である。E13 の「残るとしても重い 1 世代目だけ」は、年齢を分布にしても変わらない。
\item \textbf{1 世代目の年齢が広いと、積み上がりは半分ほどにしかならない}。既定の 1 世代目では $20\Msun$ で 0.60、$15\Msun$ で 0.47。若い 1 世代目は、2 世代目と同じように回転を保つからである。
\item \textbf{2 世代目が遅いと、窓は閉じる}。年齢が $3\times10^8\yr$ 以上なら、$20\Msun$ 以下で 0.24 以下である。
\item \textbf{AGN 円盤の 2 世代目なら、窓は広く開く}(0.67--1.00)。ただし円盤では降着でスピンが戻りうるので(E6)、この組の数値はそのままでは使えない。
\end{itemize}

\subsection{集団解析での使い方}

\paragraph{呼び方}
事象の標本ごとに、赤方偏移から年齢の節を出し、$\chi_\mathrm{SR}$ を節の上で引いて平均する。
\begin{verbatim}
T = am.nodes(gen, z, K)                             # 形は z の形 + (K,)
chi, regime = model(M[:, None], chi0[:, None], T, mu, f)
\end{verbatim}
節は等確率の分位点で、重みはどれも $1/K$ である。確率を節の平均で出すと $1/K$ 刻みになり、$K=400$ との差は $K=32$ で 0.016、$K=100$ で 0.005 だった。年齢の模型のパラメータを変えるたびに模型を作り直せばよい(0.04 秒)。

\paragraph{何が決まるか}
年齢をすべて $k$ 倍にして $f$ を $k^{-1/2}$ 倍にしても、固定点の領域のスピンは変わらない。2 世代目の年齢の範囲を 10 倍にして $f$ を $10^{-1/2}$ 倍にすると、$P_\mathrm{keep}$ の差は 0.003 以下である。したがって、データが決めるのは $f^2\times$(年齢)である。
\begin{itemize}
\item 1 世代目の年齢は、宇宙の年齢と星形成史で桁が決まる(既定で中央値 1.5 Gyr、古い組で 6.6 Gyr)。$f$ の値は、この物差しで測ることになる。物差しの不定性($\kappa_1=0$ か 1 かで中央値が 4.4 倍)は、$f$ に 2.1 倍の不定性として残る。
\item 2 世代目の年齢の尺度($T_\mathrm{lo}$ や $T_\mathrm{med}$)は、$f$ と一緒にデータから制限される量である。事前分布を狭く置くと、それがそのまま窓の有無を決める(表~\ref{tab:e22W} の行の違い)。
\item 提案: 主解析は、1 世代目を既定の組($\kappa_1=1$、$T_\mathrm{min,1}=10^7\yr$)、2 世代目を対数正規にして $T_\mathrm{med}$ を $10^5$--$10^{9.5}\yr$ の対数一様、$s$ を 0.3--1.5 dex で周辺化する。頑健さの節で、$\kappa_1=0$、星形成率の重みなし、2 世代目を対数一様にした場合を示す。事前分布の範囲は、共同研究者と決める。
\end{itemize}

\begin{result}
\begin{itemize}
\item 年齢の模型を \texttt{share/ages/ages.py} に書いた。1 世代目は $p_1\propto T^{-\kappa_1}\psi(t_m-T)$、2 世代目は $p_2\propto p_D(T)R_1(t_m-T)$ で、年齢の分位点を求積の節として返す。
\item 既定の 1 世代目の年齢は中央値 1.5 Gyr($z_m=0.2$)で、16--84\% は $5.5\times10^7$--$6.9\times10^9\yr$。これまでの 6 Gyr は $\kappa_1=0$ の組(中央値 6.6 Gyr)にあたる。
\item 窓の条件は年齢の比 $T_1/T_2\ge R(M)$ で書け、$R$ は $f$ によらない: $10$、15、20、$25\Msun$ で 299、36、7.4、1.9($\chi_b=0.78$、$\delta=0.05$)。
\item 年齢について周辺化すると、同じ $f$ で両立する度合いは、星団の目安の 2 世代目と古い 1 世代目で $20\Msun$ が 0.88、$10\Msun$ が 0.33。1 世代目を既定にすると $20\Msun$ で 0.60。2 世代目が $3\times10^8\yr$ 以上なら $20\Msun$ 以下で 0.24 以下。
\item $f$ を測る物差しは 1 世代目の年齢で、2 世代目の年齢の尺度は $f$ と一緒に制限される量である。
\end{itemize}
\end{result}

\begin{caveat}
\begin{itemize}
\item \textbf{年齢の模型は仮定である}。$\chi_\mathrm{SR}$ のような精度の約束はない。2 世代目の四つの組は目安で、文献の図から読んだ値と、増分の時間という解釈に基づく(E9 の訂正のとおり、読み方に不定性がある)。
\item \textbf{1 世代目は一つの形で書いた}。孤立連星と星団の 1 世代目の混合は、$\kappa_1$ を 0 から 1 のあいだで動かすことで近似している。金属量による形成効率の違いは入れていない(入れると、古い側に寄る)。BH ができるまでの 3--20 Myr も無視した(既定の組で $3\times10^7\yr$ より若い 10\% にだけ効く)。
\item \textbf{親の合体率の重み}には、1 世代目の全体の合体率を使った。星団の中で残った合体だけの時間変化とは違いうる。
\item \textbf{$W$、$P_\mathrm{keep}$、$P_\mathrm{drop}$ は説明用の量である}。1 世代目の生まれたスピンを 0.7 に固定し、判定の幅 $\delta=0.05$ を使っている。集団解析では、生まれたスピンの分布と測定の誤差を入れて、$\chi_\mathrm{SR}$ を測定と直接比べる(D9)。
\item 数値は $\mu=0.71\times10^{-12}\eV$、主要項の $\gamma_\mathrm{BH}$。相対論的な補正を入れると $f$ は約 1.6 倍に動くが(E12)、年齢の比で書いた条件 $R(M)$ には、1 世代目と 2 世代目で補正が同じなら効かない。
\item AGN 円盤の組は、年齢が $\chi_\mathrm{SR}$ の照合の範囲($10^5\yr$ 以上)より短い部分を含み、降着によるスピンの戻りも入れていない。
\item 2 世代目の生まれたときのスピンの分布(E10 の \texttt{hbr16.py})を親の集団とつなぐ作業は、まだである。
\end{itemize}
\end{caveat}
